\documentclass[../../main.tex]{subfiles}
\begin{document}

{\bf [解]}

\paragraph{(1)}

与えられた不等式
\begin{align}
  \dfrac{1}{2(k+1)}<\int_0^1\dfrac{1-x}{k+x}dx<\dfrac{1}{2k} \label{eq:1}
\end{align}
について考える．真ん中の項について
\begin{align}
  a_k = \int_{0}^{1} \dfrac{1-x}{x+k} dx \label{eq:2}
\end{align}
とおいて積分すると
\begin{align}
  a_k
  &= \int_{0}^{1} \left\{ -1 + \dfrac{k+1}{x+k} \right\} dx \\
  &= \left[ -x + (k+1)\log(x+k) \right]_{0}^{1} \\
  &= -1 + (k+1)\left[\log\left(k+1\right)-\log k \right] \\
  &= -1 + (k+1)\log\dfrac{k+1}{k}
\end{align}
だから $k \in \mathbb{N}$ より\cref{eq:1}に代入して
\begin{align}
  & \dfrac{1}{2(k+1)} < -1 + (k+1)\log\dfrac{k+1}{k} < \dfrac{1}{2k} \\
  \therefore
  & \dfrac{1}{2(k+1)^2}+\dfrac{1}{k+1} <  \log\dfrac{k+1}{k} < \dfrac{1}{2k}+\dfrac{1}{1+k} \label{eq:3}\\
  \therefore
  & \dfrac{2k+3}{2(k+1)^2} < \log\dfrac{k+1}{k} < \dfrac{2k+1}{2k(k+1)} \\
  \therefore
  & \dfrac{2k+3}{2(k+1)^2} < \int_{k}^{k+1} \dfrac{1}{x} dx < \dfrac{2k+1}{2k(k+1)} \label{eq:4}
\end{align}
である．

% TODO:
% \quad \left( y = \dfrac{1}{x}, y'' = \dfrac{2}{x^3} \text{ より } y = \dfrac{1}{x} \text{ は下に凸} \right)
\begin{figure}[htb]
  \centering
  \begin{tikzpicture}[x=4.0cm, y=3.0cm, >=stealth, every node/.style={font=\small}]

    % ============================================================
    % 1. 定数・厳密パラメータ
    % 例として k = 1.0 (描画上の見やすさのため)
    % ============================================================
    \def\kVal{1.0}
    \pgfmathsetmacro{\kpOne}{\kVal + 1.0}     % 2.0
    \pgfmathsetmacro{\kMid}{\kVal + 0.5}      % 1.5

    \def\C{2.0} % スケール定数 (y = C / x)
    \pgfmathsetmacro{\yL}{\C / \kVal}         % 2.0
    \pgfmathsetmacro{\yR}{\C / \kpOne}        % 1.0
    \pgfmathsetmacro{\yMid}{\C / \kMid}       % 1.3333

    % 中点での接線の傾きと切片
    \pgfmathsetmacro{\slope}{-\C / (\kMid * \kMid)} % -2.0 / 2.25 ≈ -0.8889
    \pgfmathsetmacro{\yTanL}{\slope * (\kVal - \kMid) + \yMid}   % ≈ 1.7778
    \pgfmathsetmacro{\yTanR}{\slope * (\kpOne - \kMid) + \yMid}  % ≈ 0.8889

    \coordinate (O)     at (0, 0);
    \coordinate (X_L)   at (\kVal, 0);
    \coordinate (X_M)   at (\kMid, 0);
    \coordinate (X_R)   at (\kpOne, 0);

    \coordinate (P_L)   at (\kVal, \yL);
    \coordinate (P_M)   at (\kMid, \yMid);
    \coordinate (P_R)   at (\kpOne, \yR);

    \coordinate (T_L)   at (\kVal, \yTanL);
    \coordinate (T_R)   at (\kpOne, \yTanR);

    % ============================================================
    % 2. 座標軸
    % ============================================================
    \draw[->, thick, gray!80!black] (-0.1, 0) -- (2.3, 0) node[right] {$x$};
    \draw[->, thick, gray!80!black] (0, -0.15) -- (0, 2.3) node[above] {$y$};
    \node[below left=2pt] at (O) {$O$};

    % ============================================================
    % 3. 垂線と区間境界
    % ============================================================
    \draw[thick, gray!70!black] (X_L) -- (P_L);
    \draw[thick, gray!70!black] (X_R) -- (P_R);
    \draw[dotted, gray!60] (X_M) -- (P_M);

    \node[below=2pt] at (X_L) {$k$};
    \node[below=2pt] at (X_R) {$k+1$};
    \node[below=2pt, font=\scriptsize] at (X_M) {$k+\frac{1}{2}$};

    % y 軸への投影
    \draw[dashed, gray!50] (P_L) -- (0, \yL) node[left=2pt, font=\scriptsize] {$\frac{1}{k}$};
    \draw[dashed, gray!50] (P_R) -- (0, \yR) node[left=2pt, font=\scriptsize] {$\frac{1}{k+1}$};

    % ============================================================
    % 4. 曲線 y = 1/x (下に凸)
    % ============================================================
    \draw[line width=1.5pt, blue!85!black, domain=0.85:2.2, samples=70, smooth]
    plot (\x, {\C / \x});
    \node[blue!90!black, font=\footnotesize\bfseries, above right] at (2.05, 1.05) {$y = \frac{1}{x}$};

    % ============================================================
    % 5. 弦（上側台形の上底・下底）と 接線（下側台形）
    % ============================================================
    % 弦 (緑系)
    \draw[line width=1.3pt, teal!85!black] (P_L) -- (P_R)
    node[midway, above right=1pt, font=\scriptsize\bfseries] {弦};

    % 中点での接線 (赤系)
    \draw[line width=1.2pt, red!85!black] (T_L) -- (T_R)
    node[midway, below left=1pt, font=\scriptsize\bfseries] {中点での接線};

    % 点プロット
    \fill[black] (P_L) circle (1.2pt);
    \fill[black] (P_R) circle (1.2pt);
    \fill[red!85!black] (P_M) circle (1.3pt);

  \end{tikzpicture}
  \caption{下に凸な曲線 $y = \dfrac{1}{x}$ と接線台形・弦台形による面積評価}
  \label{fig:trapezoid_log_evaluation}
\end{figure}

ここで真ん中の積分について上図のように上下から台形で挟んで面積を比較して，
\begin{align}
  & \dfrac{1}{k+1/2} < \int_{k}^{k+1} \dfrac{1}{x} dx < \dfrac{1}{2}\left(\dfrac{1}{k} + \dfrac{1}{k+1}\right) \\
  \therefore
  & \dfrac{2}{2k+1} < \log\dfrac{k+1}{k} < \dfrac{2k+1}{2k(k+1)} \label{eq:5}
\end{align}
であり，\cref{eq:4}と比べて
\begin{align}
  \dfrac{2}{2k+1} - \dfrac{2k+3}{2(k+1)^2}
  = \dfrac{1}{2(k+1)^2(2k+1)} > 0
\end{align}
とあわせて，
\begin{align}
  \dfrac{2k+3}{2(k+1)^2} < \log\dfrac{k+1}{k} < \dfrac{2k+1}{2k(k+1)}
\end{align}
だから\cref{eq:4}は示された．
よって不等式は成立である．

\paragraph{(2)}

与えられた不等式
\begin{align}
  \dfrac{m-n}{2(m+1)(n+1)}<\log\dfrac{m}{n}-\sum_{k=n+1}^m\dfrac{1}{k}<\dfrac{m-n}{2mn} \label{eq:6}
\end{align}
を示す．

\cref{eq:3}で$k=n,n+1,\cdots,m-1$として足し合わせて
\begin{align}
  \sum_{k=n}^{m-1} \dfrac{1}{k+1} + \sum_{k=n}^{m-1} \dfrac{1}{2(k+1)^2}
  < \log\dfrac{m}{n}
  < \dfrac{1}{2}\sum_{k=n}^{m-1}\left(\dfrac{1}{k} + \dfrac{1}{k+1}\right)
\end{align}
である．和の区間を$[m+1,n]$に合わせて
\begin{align}
  & \sum_{k=n+1}^{m} \dfrac{1}{k} + \sum_{k=n+1}^{m} \dfrac{1}{2k^2}
  < \log\dfrac{m}{n}
  < \sum_{k=n+1}^{m}\dfrac{1}{k} + \dfrac{1}{2}\left[\dfrac{1}{n} - \dfrac{1}{m-1}\right] \\
  \therefore
  & \sum_{k=n+1}^{m} \dfrac{1}{2k^2}
  < \log\dfrac{m}{n} - \sum_{k=n+1}^{m} \dfrac{1}{k}
  < \dfrac{m-n}{2mn} \label{eq:7}
\end{align}
である．よって\cref{eq:6}の右側の不等式は示された．

% TODO: 面積
\begin{figure}[htb]
  \centering
  \begin{tikzpicture}[x=1.8cm, y=4.0cm, >=stealth, every node/.style={font=\small}]

    % ============================================================
    % 1. 座標軸
    % ============================================================
    \draw[->, thick, gray!80!black] (-0.3, 0) -- (5.6, 0) node[right] {$x$};
    \draw[->, thick, gray!80!black] (0, -0.08) -- (0, 1.25) node[above] {$y$};
    \node[below left=2pt] at (0, 0) {$O$};

    % ============================================================
    % 2. 階段状の長方形（和 \sum 1/k^2 の領域：斜線ハッチング）
    % 区間: x \in [k, k+1] で高さ 1/k^2
    % 例示として k = 1 (n+1 に対応), k = 2, ..., k = 4 (m に対応)
    % ============================================================
    % 曲線関数 y = 1 / x^2
    % k=1: [1, 2] x [0, 1.0]
    \fill[pattern=north east lines, pattern color=blue!70, opacity=0.8]
    (1, 0) rectangle (2, 1.0);
    \draw[thick, gray!30!black] (1, 0) rectangle (2, 1.0);

    % k=2: [2, 3] x [0, 0.25]
    \fill[pattern=north east lines, pattern color=blue!70, opacity=0.8]
    (2, 0) rectangle (3, 0.25);
    \draw[thick, gray!30!black] (2, 0) rectangle (3, 0.25);

    % k=3: [3, 4] x [0, 0.111] (省略破線部)
    \node at (3.5, 0.15) {$\cdots$};

    % k=m: [4, 5] x [0, 0.0625]
    \fill[pattern=north east lines, pattern color=blue!70, opacity=0.8]
    (4, 0) rectangle (5, 0.08);
    \draw[thick, gray!30!black] (4, 0) rectangle (5, 0.08);

    % ============================================================
    % 3. 曲線 y = 1/x^2 のプロット
    % ============================================================
    \draw[line width=1.5pt, red!85!black, domain=0.95:5.3, samples=80, smooth]
    plot (\x, {1.0 / (\x * \x)});
    \node[red!90!black, font=\footnotesize\bfseries, above right] at (4.8, 0.12) {$y = \frac{1}{x^2}$};

    % ============================================================
    % 4. 積分領域との境界・目盛り
    % ============================================================
    \node[below=2pt] at (1, 0) {$n+1$};
    \node[below=2pt] at (2, 0) {$n+2$};
    \node[below=2pt] at (4, 0) {$m$};
    \node[below=2pt] at (5, 0) {$m+1$};

    % 高さのラベル
    \draw[dashed, gray!50] (1, 1.0) -- (0, 1.0)
    node[left=2pt, font=\scriptsize] {$\frac{1}{(n+1)^2}$};
    \draw[dashed, gray!50] (2, 0.25) -- (0, 0.25)
    node[left=2pt, font=\scriptsize] {$\frac{1}{(n+2)^2}$};

  \end{tikzpicture}
  \caption{級数 $\displaystyle\sum_{k=n+1}^m \dfrac{1}{k^2}$ の長方形和（斜線部）と単調減少な曲線 $y = \dfrac{1}{x^2}$ の面積比較}
  \label{fig:sum_inverse_square_comparison}
\end{figure}

左辺の$\sum_{k=n+1}^{m} \dfrac{1}{2k^2}$ は上図の斜線部の領域だから，面積を比較して
\begin{align}
  \sum_{k=n+1}^{m} \dfrac{1}{k^2}
  &> \int_{n+1}^{m+1} \dfrac{1}{x^2} dx
  = \left[ -\dfrac{1}{x} \right]_{n+1}^{m+1} \\
  &= \dfrac{1}{n+1} - \dfrac{1}{m+1}
  = \dfrac{m-n}{(m+1)(n+1)}
\end{align}
だから，\cref{eq:7}に代入して
\begin{align}
  \dfrac{m-n}{2(m+1)(n+1)}
  < \log\dfrac{m}{n} - \sum_{k=n+1}^{m} \dfrac{1}{k}
  < \dfrac{m-n}{2mn}
\end{align}
である．よって題意は示された．

{\bf [別解]}

\subparagraph{(1)}

% TODO: \dfrac{1-x}{x+k}を上下から評価した図
\begin{figure}[htb]
  \centering
  \begin{tikzpicture}[x=4.5cm, y=5.0cm, >=stealth, every node/.style={font=\small}]

    % ============================================================
    % 1. 定数・厳密パラメータ
    % 例として k = 1.0 (描画のメリハリのため)
    % g(x) = (1 - x) / (x + k)
    % ============================================================
    \def\kVal{1.0}
    \pgfmathsetmacro{\yB}{1.0 / \kVal}           % B(0, 1/k) = (0, 1.0)
    \pgfmathsetmacro{\yC}{1.0 / (\kVal + 1.0)}   % C(0, 1/(k+1)) = (0, 0.5)

    \coordinate (O) at (0, 0);
    \coordinate (A) at (1, 0);
    \coordinate (B) at (0, \yB);
    \coordinate (C) at (0, \yC);

    % ============================================================
    % 2. 面積のハッチング（曲線下の定積分領域）
    % ============================================================
    \fill[pattern=north east lines, pattern color=blue!65, opacity=0.7]
    (O) -- (B)
    -- plot[domain=0:1, samples=60, smooth] (\x, {(1.0 - \x) / (\x + \kVal)})
    -- cycle;

    % ============================================================
    % 3. 座標軸
    % ============================================================
    \draw[->, thick, gray!80!black] (-0.15, 0) -- (1.25, 0) node[right] {$x$};
    \draw[->, thick, gray!80!black] (0, -0.08) -- (0, 1.25) node[above] {$y$};
    \node[below left=2pt] at (O) {$O$};

    % ============================================================
    % 4. 評価直線（弦 AB と 接線 AC）
    % ============================================================
    % 上からの評価: 弦 AB (緑系)
    \draw[line width=1.3pt, teal!85!black] (A) -- (B);

    % 下からの評価: x=1 での接線 AC (赤系)
    % 傾き -1/(k+1)
    \draw[line width=1.3pt, red!85!black] (A) -- (C);

    % ============================================================
    % 5. 曲線 y = (1-x)/(x+k) の描画 (下に凸)
    % ============================================================
    \draw[line width=1.5pt, blue!85!black, domain=0:1, samples=80, smooth]
    plot (\x, {(1.0 - \x) / (\x + \kVal)});
    \node[blue!90!black, font=\footnotesize\bfseries, above right] at (0.35, 0.55)
    {$y = \frac{1-x}{x+k}$};

    % ============================================================
    % 6. 各点のプロットとラベル
    % ============================================================
    \fill[black] (A) circle (1.2pt) node[below=2pt] {$A(1, 0)$};
    \fill[teal!85!black] (B) circle (1.3pt) node[left=2pt] {$B\left(0, \frac{1}{k}\right)$};
    \fill[red!85!black]  (C) circle (1.3pt) node[left=2pt] {$C\left(0, \frac{1}{k+1}\right)$};

  \end{tikzpicture}
  \caption{下に凸な曲線 $y = \dfrac{1-x}{x+k}$ と2つの直角三角形 $\triangle OAC, \triangle OAB$ による面積評価}
  \label{fig:alternative_integral_triangle_evaluation}
\end{figure}

$f(x)=\dfrac{1-x}{x+k}$とおくと，
\begin{align}
  f(x) =  -1 + \dfrac{k+1}{x+k}
\end{align}
でこれは下に凸だから，$y=f(x)$のグラフの概形は上図のようになる．
下からは$x=1$での接線で，上からは$(0,1/k),(1,0)$を結んだ線分で挟むと，
\begin{align}
  \triangle OAC < \int_{0}^{1}f(x)dx < \triangle OAB
\end{align}
である．従って
\begin{align}
  \dfrac{1}{2(k+1)} < \int_{0}^{1}f(x)dx < \dfrac{1}{2k}
\end{align}
を得る．

\subparagraph{(2)}
\cref{eq:7}から左側の不等式は以下のように積分なしでも全く同じ評価ができる．
\begin{align}
  \sum_{k=n+1}^{m} \dfrac{1}{k^2} > \sum_{k=n+1}^{m} \dfrac{1}{k(k+1)} = \dfrac{1}{n+1} - \dfrac{1}{m+1}
\end{align}

\end{document}
